
\chapter{2017年上海卷（秋）}

\section{填空题}

\begin{problem}
已知集合 \(A=\{1,2,3,4\}\)，集合 \(B=\{3,4,5\}\)，则 \(A\cap B=\fillinblank{}\).
\end{problem}

\begin{problem}
若排列数 \(\mathrm A_6^m=6\times5\times4\)，则 \(m=\fillinblank{}\).
\end{problem}

\begin{problem}
不等式 \(\dfrac{x-1}{x}>1\) 的解集为\fillinblank{}.
\end{problem}

\begin{problem}
已知球的体积为 \(36\pi\)，则该球主视图的面积等于\fillinblank{}.
\end{problem}

\begin{problem}
已知复数 \(z\) 满足 \(z+\dfrac3z=0\)，则 \(\lvert z\rvert=\fillinblank{}\).
\end{problem}

\begin{problem}
设双曲线 \(\dfrac{x^2}{9}-\dfrac{y^2}{b^2}=1\)（\(b>0\)）的焦点为 \(F_1\)，\(F_2\)，\(P\) 为该双曲线上的一点，若 \(\lvert PF_1\rvert=5\)，则 \(\lvert PF_2\rvert=\fillinblank{}\).
\end{problem}

\begin{problem}
如图，以长方体 \(ABCD-A_1B_1C_1D_1\) 的顶点 \(D\) 为坐标原点，过 \(D\) 的三条棱所在的直线为坐标轴，建立空间直角坐标系. 若 \(\overrightarrow{DB_1}\) 的坐标为 \((4,3,2)\)，则 \(\overrightarrow{AC_1}\) 的坐标是\fillinblank{}.

\bitmapfigure[max width=6.0cm]{img/2017/shanghai/q7_fig1.png}
\end{problem}

\begin{problem}
定义在 \((0,+\infty)\) 上的函数 \(y=f(x)\) 的反函数为 \(y=f^{-1}(x)\). 若
\(g(x)=\begin{cases}
3^x-1,&x\le 0,\\
f(x),&x>0,
\end{cases}\)
且 \(g(x)\) 为奇函数，则 \(f^{-1}(x)=2\) 的解为\fillinblank{}.
\end{problem}

\begin{problem}
已知四个函数：
\begin{circlelist}
\item \(y=-x\)；
\item \(y=-\dfrac1x\)；
\item \(y=x^3\)；
\item \(y=x^{\frac12}\).
\end{circlelist}
从中任选 \(2\) 个，则事件“所选 \(2\) 个函数的图像有且仅有一个公共点”的概率为\fillinblank{}.
\end{problem}

\begin{problem}
已知数列 \(\{a_n\}\) 和 \(\{b_n\}\)，其中 \(a_n=n^2\)，\(n\in\N^*\)，\(\{b_n\}\) 的项是互不相等的正整数. 若对于任意 \(n\in\N^*\)，\(\{b_n\}\) 的第 \(a_n\) 项都等于 \(\{a_n\}\) 的第 \(b_n\) 项，则 \(\dfrac{\lg(b_1b_4b_9b_{16})}{\lg(b_1b_2b_3b_4)}=\fillinblank{}\).
\end{problem}

\begin{problem}
设 \(\alpha_1\)，\(\alpha_2\in\R\)，且 \(\dfrac1{2+\sin\alpha_1}+\dfrac1{2+\sin(2\alpha_2)}=2\)，则 \(\lvert10\pi-\alpha_1-\alpha_2\rvert\) 的最小值等于\fillinblank{}.
\end{problem}

\begin{problem}
如图，用 \(35\) 个单位正方形拼成一个矩形，点 \(P_1\)，\(P_2\)，\(P_3\)，\(P_4\) 以及四个标记为“▲”的点在正方形的顶点处，设集合 \(\Omega=\{P_1,P_2,P_3,P_4\}\)，点 \(P\in\Omega\)，过 \(P\) 作直线 \(l_P\)，使得不在 \(l_P\) 上的“▲”的点分布在 \(l_P\) 的两侧. 用 \(D_1(l_P)\) 和 \(D_2(l_P)\) 分别表示 \(l_P\) 一侧和另一侧的“▲”的点到 \(l_P\) 的距离之和. 若过 \(P\) 的直线 \(l_P\) 中有且只有一条满足 \(D_1(l_P)=D_2(l_P)\)，则 \(\Omega\) 中所有这样的 \(P\) 为\fillinblank{}.

\bitmapfigure[max width=6.0cm]{img/2017/shanghai/q12_fig1.png}
\end{problem}

\section{单选题}

\begin{problem}
关于 \(x\)，\(y\) 的二元一次方程组
\(\begin{cases}
x+5y=0,\\
2x+3y=4
\end{cases}\)的系数行列式 \(D\) 为（\quad）

\choices
    {\(\left|\begin{matrix}0&5\\ 4&3\end{matrix}\right|\)}
    {\(\left|\begin{matrix}1&0\\ 2&4\end{matrix}\right|\)}
    {\(\left|\begin{matrix}1&5\\ 2&3\end{matrix}\right|\)}
    {\(\left|\begin{matrix}6&0\\ 5&4\end{matrix}\right|\)}
\end{problem}

\begin{problem}
在数列 \(\{a_n\}\) 中，\(a_n=\left(-\dfrac12\right)^n\)，\(n\in\N^*\)，则 \(\displaystyle\lim_{n\to\infty}a_n\)（\quad）

\choices
    {\(-\dfrac12\)}
    {\(0\)}
    {\(\dfrac12\)}
    {\text{不存在}}
\end{problem}

\begin{problem}
已知 \(a\)，\(b\)，\(c\) 为实常数，数列 \(\{x_n\}\) 的通项 \(x_n=an^2+bn+c\)，\(n\in\N^*\)，则“存在 \(k\in\N^*\)，使得 \(x_{100+k}\)，\(x_{200+k}\)，\(x_{300+k}\) 成等差数列”的一个必要条件是（\quad）

\choices
    {\(a\ge0\)}
    {\(b\le0\)}
    {\(c=0\)}
    {\(a-2b+c=0\)}
\end{problem}

\begin{problem}
在平面直角坐标系 \(xOy\) 中，已知椭圆 \(C_1:\dfrac{x^2}{36}+\dfrac{y^2}{4}=1\) 和 \(C_2:x^2+\dfrac{y^2}{9}=1\). \(P\) 为 \(C_1\) 上的动点，\(Q\) 为 \(C_2\) 上的动点，\(w\) 是 \(\overrightarrow{OP}\cdot\overrightarrow{OQ}\) 的最大值. 记 \(\Omega=\{(P,Q)\mid P\text{ 在 }C_1\text{ 上，}Q\text{ 在 }C_2\text{ 上，且 }\overrightarrow{OP}\cdot\overrightarrow{OQ}=w\}\)，则 \(\Omega\) 中元素个数为（\quad）

\choices
    {\(2\)个}
    {\(4\)个}
    {\(8\)个}
    {\text{含有无穷个元素}}
\end{problem}

\section{解答题}

\begin{problem}
如图，直三棱柱 \(ABC-A_1B_1C_1\) 的底面为直角三角形，两直角边 \(AB\) 和 \(AC\) 的长分别为 \(4\) 和 \(2\)，侧棱 \(AA_1\) 的长为 \(5\).

\begin{enumerate}
\item 求三棱柱 \(ABC-A_1B_1C_1\) 的体积；
\item 设 \(M\) 是 \(BC\) 中点，求直线 \(A_1M\) 与平面 \(ABC\) 所成角的大小.
\end{enumerate}

\bitmapinlinefigure[0.9\linewidth][max width=6.0cm]{img/2017/shanghai/q17_fig1.png}
\end{problem}

\begin{problem}
已知函数 \(f(x)=\cos^2x-\sin^2x+\dfrac12\)，\(x\in(0,\pi)\).
\begin{enumerate}
\item 求 \(f(x)\) 的单调递增区间；
\item 设 \(\triangle ABC\) 为锐角三角形，角 \(A\) 所对边 \(a=\sqrt{19}\)，角 \(B\) 所对边 \(b=5\). 若 \(f(A)=0\)，求 \(\triangle ABC\) 的面积.
\end{enumerate}
\end{problem}

\begin{problem}
根据预测，某地第 \(n\)（\(n\in\N^*\)）个月共享单车的投放量和损失量分别为 \(a_n\) 和 \(b_n\)（单位：辆），其中
\(a_n=\begin{cases}
5n^4+15,&1\le n\le 3,\\
-10n+470,&n\ge4,
\end{cases}\quad b_n=n+5\)，第 \(n\) 个月底的共享单车的保有量是前 \(n\) 个月的累计投放量与累计损失量的差.

\begin{enumerate}
\item 求该地区第 \(4\) 个月底的共享单车的保有量；
\item 已知该地共享单车停放点第 \(n\) 个月底的单车容纳量 \(S_n=-4(n-46)^2+8800\)（单位：辆）. 设在某月底，共享单车保有量达到最大，问该保有量是否超出了此时停放点的单车容纳量？
\end{enumerate}
\end{problem}

\begin{problem}
在平面直角坐标系 \(xOy\) 中，已知椭圆 \(\Gamma:\dfrac{x^2}{4}+y^2=1\). \(A\) 为 \(\Gamma\) 的上顶点，\(P\) 为 \(\Gamma\) 上异于上、下顶点的动点，\(M\) 为 \(x\) 轴正半轴上的动点.

\begin{enumerate}
\item 若 \(P\) 在第一象限，且 \(\lvert OP\rvert=\sqrt2\)，求 \(P\) 的坐标；
\item 设 \(P\left(\dfrac85,\dfrac35\right)\)，若以 \(A\)，\(P\)，\(M\) 为顶点的三角形是直角三角形，求 \(M\) 的横坐标；
\item 若 \(\lvert MA\rvert=\lvert MP\rvert\)，直线 \(AQ\) 与 \(\Gamma\) 交于另一点 \(C\)，且 \(\overrightarrow{AQ}=2\overrightarrow{AC}\)，\(\overrightarrow{PQ}=4\overrightarrow{PM}\)，求直线 \(AQ\) 的方程.
\end{enumerate}
\end{problem}

\begin{problem}
设定义在 \(\R\) 上的函数 \(f(x)\) 满足：对于任意的 \(x_1\)，\(x_2\in\R\)，当 \(x_1<x_2\) 时，都有 \(f(x_1)\le f(x_2)\).

\begin{enumerate}
\item 若 \(f(x)=ax^3+1\)，求 \(a\) 的取值范围；
\item 若 \(f(x)\) 为周期函数，证明：\(f(x)\) 是常值函数；
\item 设 \(f(x)\) 恒大于零，\(g(x)\) 是定义在 \(\R\) 上、恒大于零的周期函数，\(M\) 是 \(g(x)\) 的最大值. 函数 \(h(x)=f(x)g(x)\). 证明：“\(h(x)\) 是周期函数”的充要条件是“\(f(x)\) 是常值函数”.
\end{enumerate}
\end{problem}
